% SEGMENT OLP-0487-S01
% Part: applied-modal-logics
% Chapter: epistemic-logic
% Section: properties-accessibility

\documentclass[../../../include/open-logic-section]{subfiles}

\begin{document}

\olfileid{aml}{el}{acc}

\olsection{அணுகுறவுகளும் அறிவுநிலைக் கொள்கைகளும்}

இயல்பான வாய்ப்புநிலைத் தருக்கங்களின் சட்டகப் பொருத்தத்தைப் பற்றி நாம்
ஏற்கெனவே அறிந்தவற்றை வைத்துக்கொண்டு, அறிவுநிலைப் பொருள்கோள்களின் கீழ்
சிறப்பியல்பு !!{formula}s எவ்வாறு இருக்கும் என்று காண விரும்பலாம்.
அறிவுநிலைத் தருக்கங்கள் பொதுவாக S5 இல் பொருள்கொள்ளப்படுகின்றன என்று
ஏற்கெனவே கூறியுள்ளோம். எனவே பல்வேறு அறிவுநிலைக் கொள்கைகள் எவ்வாறு
குறிக்கப்படுகின்றன என்றும் அவை பல்வேறு சட்டக நிபந்தனைகளுடன் எவ்வாறு
பொருந்துகின்றன என்றும் காண்போம்.

வெவ்வேறு வாய்ப்புநிலை !!{formula}s, அணுகுறவுகளின் வெவ்வேறு பண்புகளைச்
சிறப்பித்தன என்பதை இயல்பான வாய்ப்புநிலைத் தருக்கத்திலிருந்து நினைவுகூர்க.
குறிப்பிட்ட அறிவுநிலைக் கொள்கைகளுக்குப் பொருந்தும் சிலவற்றை இவ்வட்டவணை
தனியே காட்டுகிறது.

\begin{table}[t]
    \begin{tabular}{| p{.48\textwidth} || p{.48\textwidth} |}
      \hline
      {\emph{$R$ \dots எனில்}} & {\emph{\dots ஆனது~$\mModel{M}$ இல் மெய்யாகும்:}} \\
      \hline \hline

      & $\Knows (p \lif q) \lif (\Knows p \lif \Knows q)$
      \hfill \newline{(அடைவு)} \\
      \hline
      \emph{தற்சுட்டுப் பண்புள்ளது}: $\forall w Rww$
      & $\Knows p \lif p$ \hfill \newline{(மெய்மைத்தன்மை)} \\
      \hline
      \emph{கடப்புப் பண்புள்ளது}: \newline
      $\forall u \forall v \forall w ((Ruv \land Rvw) \lif Ruw)$ &
      $\Knows p \lif \Knows \Knows p$ \hfill
           \newline{(நேர்முகத் தன்னோக்கு)} \\
      \hline
      \emph{யூக்ளிடியப் பண்புள்ளது}: \newline
      $\forall w \forall u \forall v ((Rwu \land Rwv) \lif Ruv)$ &
      $\lnot \Knows p \lif \Knows \neg \Knows p$ \hfill
        \newline{(எதிர்மறைத் தன்னோக்கு)} \\
      \hline
    \end{tabular}
    \caption{நான்கு அறிவுநிலைக் கொள்கைகள்.}
    \ollabel{tab:four}
  \end{table}

$T$ கோட்பாட்டுறுதிக்குப் பொருந்தும் மெய்மைத்தன்மை, இக்கொள்கைகளில் மிகக்
குறைவாகச் சர்ச்சைக்குரியது என்று அடிக்கடி கருதப்படுகிறது; ஏனெனில் ஓர்
!!a{formula} அறியப்பட்டால் அது மெய்யாக இருக்க வேண்டும் என்ற கூற்றை அது
குறிக்கிறது. அடைவும் நேர்முக, எதிர்மறைத் தன்னோக்குகளும் இன்னும் அதிகமாகச்
சர்ச்சைக்குரியவை.

$K$ கோட்பாட்டுறுதிக்குப் பொருந்தும் அடைவு, ஒரு முகவரின் அறிவு
நிபந்தனைவாக்கத்தின் கீழ் மூடப்பட்டுள்ளது என்ற கருத்தைக் குறிக்கிறது. சில
நேர்வுகளில் இது நமக்கு ஏற்புடையதாகத் தோன்றலாம். எடுத்துக்காட்டாக, நான்
விக்டோரியாவில் இருந்தால் வான்கூவர் தீவில் இருப்பேன் என்பதை நான் அறியலாம்.
வழக்கத்திற்கு மாறான ஐயுறவுச் சூழ்நிலைகளை ஒதுக்கினால், நான் விக்டோரியாவில்
இருப்பதை உண்மையிலேயே அறிவேன்; ஆகவே நான் வான்கூவர் தீவில் இருப்பதையும்
அறிவேன் என்பது இதனால் தெரிய வேண்டும். எனவே இந்நேர்வில், எனது அறிவின்
தருக்க அடைவு ஒப்பீட்டளவில் உள்ளுணர்வுக்கு ஏற்புடையதாகத் தோன்றலாம். மறுபுறம்,
நமது அறிவின் விளைவுகள் அனைத்தையும் நாம் எப்போதும் சிந்தித்துப் பார்ப்பதில்லை;
ஆகவே பிற நேர்வுகளில் இது உள்ளுணர்வுக்கு ஏற்பில்லாத முடிவுகளைத் தரலாம்.

சில நேரங்களில் KK-கொள்கை எனப்படும் நேர்முகத் தன்னோக்கு, நான் ஏதோவொன்றை
அறிந்தால், அதை நான் அறிவதை அறிகிறேன் என்ற கூற்றாகச் சில நேரங்களில்
வெளிப்படுத்தப்படுகிறது. இது கோட்பாட்டுறுதி~4 இன் அறிவுநிலை இணையாகும்.
அதற்கேற்ப, எதிர்மறைத் தன்னோக்கு என்பது நான் ஏதோவொன்றை \emph{அறியவில்லை}
எனில், அதை நான் அறியவில்லை என்பதை அறிகிறேன் என்ற கூற்றாக
வெளிப்படுத்தப்படுகிறது; இது கோட்பாட்டுறுதி~5 இன் இணையாகும். நான் உண்மையில்
அறியும் ஒன்றை அறிகிறேனா இல்லையா என்பது எனக்கே உறுதியாக இல்லாத, மிகவும்
சாதாரணமான எதிரெடுத்துக்காட்டுகளை இவ்விரு கொள்கைகளுக்கும் தர முடியும் என்று
தோன்றுகிறது.


\end{document}
